\documentclass[10pt,a4paper]{amsart}
\usepackage[margin=19mm]{geometry}
\usepackage{amsmath,amssymb,mathtools}
\usepackage[scr=rsfs,cal=euler]{mathalfa}
\usepackage{xfrac}
\usepackage[unicode,hidelinks,bookmarks=false]{hyperref}
\newcommand{\ie}{i.e.\ }
\newcommand{\cf}{cf.\ }
\newcommand{\resp}{respectively\ }
\newcommand{\Subsection}{\subsection}
\providecommand{\nobreakdash}{\nobreak\hbox{-}}
\setlength{\emergencystretch}{2em}
\multlinegap0pt
\usepackage{bm}
\usepackage{bbm}
\usepackage{dsfont}
\usepackage{xspace}
\usepackage{cancel}
\usepackage{mathtools}
\usepackage{amsthm}
\makeatletter
\@ifundefined{definition}{\newtheorem{definition}{Definition}}{}
\@ifundefined{proposition}{\newtheorem{proposition}[definition]{Proposition}}{}
\makeatother
\def \N {\mathbb{N}}
\def \R {\mathbb{R}}
\renewcommand{\emptyset}{\varnothing}
\renewcommand{\ge}{\geqslant}
\renewcommand{\le}{\leqslant}
\title[Nonlocal eigenvalue problems and superposition operators]{Nonlocal eigenvalue problems and superposition operators}
\author{Serena Dipierro, Edoardo Proietti Lippi, Caterina Sportelli, Enrico Valdinoci}
\date{Source: arXiv:2602.18035v1 (CC BY 4.0)}
\begin{document}
\maketitle

\noindent\emph{Source and licence.} Nonlocal eigenvalue problems and superposition operators. Version arXiv:2602.18035v1 (CC BY 4.0), distributed under the Creative Commons Attribution 4.0 International licence, as recorded in the arXiv licence metadata for this exact version and shown on the abstract page https://arxiv.org/abs/2602.18035v1. Full rights notice: licenses/arxiv-2602.18035-CC-BY-4.0.txt.\par\smallskip

\noindent\textit{Editorial scope.} Counted unit: the Proposition whose source label is \texttt{PROPD1} in the article (source0.tex lines 1742--1765, source label \texttt{PROPD1}), delivered together with the article's own complete original proof of it (source0.tex lines 1767--1851, one \texttt{proof} environment of 85 lines). No mathematics is written, repaired or re-derived: statement and proof are the article's own text line for line. Required context retained uncounted: none. Every definition, display and label of those lines is retained unchanged. Omitted from this package and not redistributed: 1--1741, 1766--1766, 1852--2292. Nothing inside a retained excerpt is omitted or altered: every retained body is delivered exactly as the article's lines are written, so no omission receipt of any kind accompanies this package. Citations: the retained excerpts carry no citation command at all, so no bibliography entry of the article is transcribed and no reference is redistributed. Reference adaptation: none was needed and none was made; every reference inside the delivered unit resolves inside it, so no unresolved reference or citation is produced. Printed slips kept literal: no printed word, symbol, hypothesis or number of the counted statement or of its proof was altered. Layout and class: the article's own class is \texttt{amsart} with packages including amsmath, amsthm, mathtools, amsfonts, amsrefs, graphicx, framed, amssymb, esvect, xcolor, eucal, mathrsfs, hyperref, babel; the wrapper is the lane's pinned \texttt{amsart} editorial wrapper at 10pt on A4, which restates the theorem-like environments the retained text needs and adds the article's own macro definitions that the retained text needs (\texttt{\textbackslash N}, \texttt{\textbackslash R}, \texttt{\textbackslash emptyset}, \texttt{\textbackslash ge}, \texttt{\textbackslash le}), with extra wrapper packages (bm, bbm, dsfont, xspace, cancel, mathtools). The delivered numbering is the unit's own. Assets: no retained span contains a figure environment, a graphics-inclusion command, a PDF-page inclusion, a table or a TikZ picture; no image file is copied into this package, the package declares no asset dependency and no image file is redistributed with it. Licence: version arXiv:2602.18035v1 of this article is distributed under the Creative Commons Attribution 4.0 International licence, as recorded in the arXiv licence metadata for this exact version and shown on the abstract page https://arxiv.org/abs/2602.18035v1. A full rights notice is delivered at licenses/arxiv-2602.18035-CC-BY-4.0.txt. Characters: every retained excerpt is pure ASCII, so the delivered engine, XeLaTeX with its default Latin Modern text font, reports no missing character.
\begin{proposition}\label{PROPD1}
Let~$M\in\N\cap[1,+\infty)$.
Let~$\Omega=\Omega_1\cup\dots\cup\Omega_M$, where, for all~$j\in\{1,\dots,M\}$, the set~$\Omega_j$ is open, bounded, connected,
with boundary of class~$C^1$,
and such that~$\overline{\Omega_i}\cap\overline{ \Omega_j}=\emptyset$ if $i\neq j$. 

Suppose that
\begin{equation}\label{lammegmbmejp}
\lambda(\Omega_1)\le\lambda(\Omega_2)\le\dots\le\lambda(\Omega_M).\end{equation}
Then:
\begin{itemize}
\item[(i).] $\lambda(\Omega)=\min\{\lambda(\Omega_1),\dots,\lambda(\Omega_M)\}=\lambda(\Omega_1)$.
\item[(ii).] If~$u_\Omega$ is any eigenfunction corresponding to the eigenvalue~$\lambda(\Omega)$,
for every~$j\in\{1,\dots,M\}$ we have that either~$u_\Omega\equiv0$ in~$\Omega_j$
or~$\lambda(\Omega)=\lambda(\Omega_j)$.
\item[(iii).] If~$j_\star\in\{1,\dots,M\}$ is such that~$\lambda(\Omega)=\lambda(\Omega_{j_\star})$
and~$u_{\Omega_{j_\star}}$ is any eigenfunction corresponding to the eigenvalue~$\lambda(\Omega_{j_\star})$,
then~$u_{\Omega_{j_\star}}$, extended to zero outside~${\Omega_{j_\star}}$,
is also an eigenfunction in~$\Omega$ corresponding to the eigenvalue~$\lambda(\Omega)$.
\item[(iv).] The eigenvalue~$\lambda(\Omega)$ is simple if and only if~$\lambda(\Omega_1)<\lambda(\Omega_2)$.
\item[(v).] If~$\lambda(\Omega_1)=\lambda(\Omega_2)<\lambda(\Omega_j)$
for all~$j\in\{3,\dots,M\}$, then~$\lambda(\Omega)$ has multiplicity two and its eigenspace is the linear combination of the eigenspaces of~$\lambda(\Omega_1)$ and~$\lambda(\Omega_2)$.
\end{itemize}
\end{proposition}

\begin{proof} 
We observe that, given~$j\in\{1,\dots,M\}$,
\begin{equation}\label{AiksptD1}\begin{split}
& \lambda(\Omega)=\min_{u\in H^1_0(\Omega)\setminus\{0\}}\frac{\displaystyle\int_\Omega|\nabla u(x)|^2\,dx}{\|u\|^2_{L^2(\Omega)}}
\le\min_{{u\in H^1_0(\Omega)\setminus\{0\}}\atop{\tiny{\mbox{$u\equiv0$ in~$\R^N\setminus\Omega_j$}}}}
\frac{\displaystyle\int_\Omega|\nabla u(x)|^2\,dx}{\|u\|^2_{L^2(\Omega)}}\\&\qquad\qquad\qquad=
\min_{u\in H^1_0(\Omega_j)\setminus\{0\}}\frac{\displaystyle\int_\Omega|\nabla u(x)|^2\,dx}{\|u\|^2_{L^2(\Omega)}}
=\lambda(\Omega_j).
\end{split}\end{equation}

Also
if~$u_\Omega$ is any eigenfunction corresponding to the eigenvalue~$\lambda(\Omega)$
and~$\chi_{\Omega_j}$ is the indicator function of~$\Omega_j$,
that is
$$ \chi_{\Omega_j}(x):=\begin{cases}1&{\mbox{ if }}x\in\Omega_j,\\
0&{\mbox{ otherwise, }}\end{cases}$$
we have that~$u_\Omega\chi_{\Omega_j}\in H^1_0(\Omega_j)$ and that
$$-\Delta(u_\Omega\chi_{\Omega_j})=-\Delta u_\Omega=\lambda(\Omega)\,u_\Omega
=\lambda(\Omega)\,u_\Omega\chi_{\Omega_j},$$
showing that \begin{equation}\label{AB7}\begin{split}&{\mbox{either~$u_\Omega\chi_{\Omega_j}$ is an eigenfunction in~$\Omega_j$ corresponding to
the eigenvalue~$\lambda(\Omega)$}}\\&{\mbox{or~$u_\Omega\chi_{\Omega_j}$ vanishes identically.}}\end{split}\end{equation}

As a consequence, 
\begin{equation}\label{AiksptD}
{\mbox{either~$\lambda(\Omega)\ge\lambda(\Omega_j)$
or~$u_\Omega$ vanishes identically in~$\Omega_j$.}}\end{equation}

Thus, pick a point~$p\in\Omega$ such that~$u_\Omega(p)\ne0$.
Let~$j_p\in\{1,\dots,M\}$ be such that~$p\in\Omega_{j_p}$. We deduce from~\eqref{AiksptD}
and~\eqref{lammegmbmejp}
that~$\lambda(\Omega)\ge\lambda(\Omega_{j_p})\ge\lambda(\Omega_1)$.
This and~\eqref{AiksptD1} (used with~$j:=1$) give that~$\lambda(\Omega)=\lambda(\Omega_1)$,
which completes the proof of~(i).

Moreover, by combining~\eqref{AiksptD1} and~\eqref{AiksptD}, the claim in~(ii) plainly follows.

Regarding the claim in~(iii), the assumptions in~(iii) yield that the inequality in~\eqref{AiksptD1}
is actually an equality, whence any function attaining the minimum
in the Rayleigh quotient corresponding to~$\lambda(\Omega_j)$
also attains the minimum
in the Rayleigh quotient corresponding to~$\lambda(\Omega)$
(up to extending this function as zero outside of~$\Omega_j$).
The proof of~(iii) is thereby complete.

To prove~(iv), 
let us first suppose that~$\lambda(\Omega)$ is simple and, for the sake of
contradiction, that~$\lambda(\Omega_1)=\lambda(\Omega_2)$.
Then, by~(i), we have that~$\lambda(\Omega)=\lambda(\Omega_{j_\star})$, for all~$j_\star\in\{1,2\}$,
and then, by~(iii), if~$
u_{\Omega_{j_\star}}$ is any eigenfunction corresponding to the eigenvalue~$\lambda(\Omega_{j_\star})$,
then~$u_{\Omega_{j_\star}}\chi_{\Omega_{j_\star}}$ is
an eigenfunction in~$\Omega$ corresponding to the eigenvalue~$\lambda(\Omega)$.

The assumed simplicity of~$\lambda(\Omega)$ thus entails that there exists~$c\in\R$ such that~$ u_{\Omega_{1}}\chi_{\Omega_{1}}=c\,u_{\Omega_{2}}\chi_{\Omega_{2}}$. But this gives that~$u_{\Omega_1}$ vanishes identically
in~$\Omega_1$, which provides the desired contradiction.

To complete the proof of~(iv), let us now suppose that~$\lambda(\Omega_1)<\lambda(\Omega_2)$.
Then, by~(i) and~\eqref{lammegmbmejp}, we have that~$\lambda(\Omega)=\lambda(\Omega_1)<\lambda(\Omega_j)$
for all~$j\in\{2,\dots,M\}$.

Hence, by~(ii), any eigenfunction in~$\Omega$ corresponding to the eigenvalue~$\lambda(\Omega)$
must vanish identically in~$\Omega_j$ for all~$j\in\{2,\dots,M\}$.
Accordingly, any eigenfunction in~$\Omega$ corresponding to the eigenvalue~$\lambda(\Omega)$
must lie in~$H^1_0(\Omega_1)$
and therefore be an eigenfunction in~$\Omega_1$ corresponding to the eigenvalue~$\lambda(\Omega)$.

Since~$\Omega_1$ is connected, we know that such an eigenfunction is unique,
up to scalar multiplication. We have thereby proved that the space of eigenfunctions
in~$\Omega$ corresponding to~$\lambda(\Omega)$ is one-dimensional, thus completing the proof of~(iv).

To prove~(v), we argue as follows. By~(iii), we know that the multiplicity of~$\lambda(\Omega)$ is at least two and that the eigenfunctions~$u_{\Omega_1}$
and~$u_{\Omega_2}$ corresponding to~$\lambda(\Omega_1)$ and~$\lambda(\Omega_2)$ (in~$\Omega_1$ and~$\Omega_2$
respectively, and extended to zero outside their domains) are also eigenfunctions
corresponding to~$\lambda(\Omega)$
(and notice that~$u_{\Omega_1}$ and~$u_{\Omega_2}$ are unique up to scalar multiplication). Hence, to complete the proof of~(v), we pick any eigenfunction~$u_\Omega$ corresponding to~$\lambda(\Omega)$ and we need to show that
\begin{equation} \label{oqwdjfrptjnmuy43tghbnl1}
{\mbox{$u_{\Omega}$ is a linear combination of~$u_{\Omega_1}$ and~$u_{\Omega_2}$. }}\end{equation}

For this objective, we use~(ii) to see that~$u_\Omega\equiv0$ outside~$\Omega_1\cup\Omega_2$ and therefore
\begin{equation} \label{oqwdjfrptjnmuy43tghbnl}
u_\Omega=u_\Omega\chi_{\Omega_1}+u_\Omega\chi_{\Omega_2}.\end{equation} 

Also, by~\eqref{AB7}, we know that either~$u_\Omega\chi_{\Omega_1}$ 
vanishes identically or it is an eigenfunction associated with~$\lambda(\Omega_1)$: either way, since~$\lambda(\Omega_1)$ is simple, there exists~$c_1\in\R$ such that~$u_\Omega\chi_{\Omega_1}=c_1u_{\Omega_1}$. In the same vein, there exists~$c_2\in\R$ such that~$u_\Omega\chi_{\Omega_2}=c_1u_{\Omega_2}$. Plugging these bits of information into~\eqref{oqwdjfrptjnmuy43tghbnl}, the desired result in~\eqref{oqwdjfrptjnmuy43tghbnl1} follows.
\end{proof}

\end{document}
